\documentclass[12pt, a4paper]{article}
\usepackage[utf8]{inputenc}
\usepackage[T1]{fontenc}
\usepackage{amsmath, amssymb, amsthm}
\usepackage{geometry}
\geometry{margin=2.5cm}
\usepackage{hyperref}
\usepackage{booktabs}
\usepackage{listings}
\usepackage{xeCJK}
\setCJKmainfont{cwTeX Q Ming}

\lstset{basicstyle=\ttfamily\small, breaklines=true, frame=single}

\title{動手玩 Factor Models：從 PCA 到 S\&P 500 的潛在因子}
\author{Sheng-Yuan Chen}
\date{2026-07-06}

\begin{document}
\maketitle

\begin{center}
完整程式碼（R 原版 + Python 重寫）：\url{https://github.com/Mullerchen910912/factor-models-in-r}
\end{center}

\section{前言：少數幾股力量，驅動一大票序列}

Factor model 的核心：很多觀察到的序列，其實被\textbf{少數幾個看不到的共同因子}驅動。這篇用 Python 對 S\&P 500 的日報酬動手做一次，讓資料自己回答：要幾個因子？因子又是什麼？

\section{一、資料整理：從價格到報酬矩陣}

\begin{lstlisting}[language=Python]
df = pd.read_csv("sp500_prices.csv", usecols=["symbol", "date", "adjusted"])
df["date"] = pd.to_datetime(df["date"])
df = df.sort_values(["symbol", "date"])
df["ret"] = df.groupby("symbol")["adjusted"].pct_change()
wide = df.pivot(index="date", columns="symbol", values="ret").dropna()
# 2,384 天 × 89 檔
\end{lstlisting}

\section{二、要幾個因子？用 PCA 看}

\begin{lstlisting}[language=Python]
X = StandardScaler().fit_transform(wide.values)
pca = PCA().fit(X)
evr = pca.explained_variance_ratio_
print(evr[0], evr[:3].sum(), evr[:5].sum())
\end{lstlisting}

\textbf{PC1 就解釋了 40.8\% 的變異}（大盤因子），前 3 個 50.8\%、前 5 個 55.2\%。少數幾個因子就抓住 89 檔股票共同起伏的一大半。

\section{三、這些因子是什麼？用 varimax + 產業標籤}

3 因子 varimax，再把載荷最高的股票貼上產業別：

\begin{center}
\begin{tabular}{lll}
\toprule
因子 & 載荷最高的股票 & 資料標出的產業 \\
\midrule
Factor 1 & ADBE, SNPS, MSFT, NVDA & Information Technology \\
Factor 2 & COP, CVX, XOM, EOG & Energy \\
Factor 3 & DUK, D, PEP, NEE, PG & Utilities／Consumer Staples \\
\bottomrule
\end{tabular}
\end{center}

我們沒告訴模型任何產業資訊，它卻自己抽出「科技、能源、防禦」三股力量。

\section{四、往因果推論走一步}

同樣的因子結構也是現代因果推論的基礎。Bai \& Wang (2024) 的因果因子模型，就是把「用因子吸收共同衝擊」用來估處理效果 —— 和 IFE／CCE panel 估計、policy targeting 是同一條線。

\section{小結}

\begin{center}
\begin{tabular}{lll}
\toprule
步驟 & 工具 & 回答的問題 \\
\midrule
資料整理 & pandas & 價格 → 報酬寬矩陣 \\
要幾個因子 & PCA & 共同變異的維度 \\
因子是什麼 & FactorAnalysis + varimax & 因子的經濟意義 \\
往因果走 & Bai-Wang 因果因子模型 & 用因子估處理效果 \\
\bottomrule
\end{tabular}
\end{center}

完整程式碼在 GitHub：\url{https://github.com/Mullerchen910912/factor-models-in-r}。

\end{document}
